% ==============================================================================
% SECTION 1: INTRODUCTION (01_intro.tex)
% VAI TRÒ PHỤ TRÁCH: Report Writer (RW)
% THỨ TỰ VIẾT: BƯỚC 3 (Sau khi đã xong §3 và §2 để nắm chắc Gap & Method)
% NGUYÊN LIỆU ĐỐI SOÁT: team-synthesis/proposal.md (§2), team-synthesis/gap-list.md
% ĐỘ DÀI MỤC TIÊU: ~1 trang, cấu trúc chuẩn 5 đoạn văn độc lập
% ==============================================================================
%
% [HƯỚNG DẪN DÀN Ý 5 ĐOẠN DÀNH CHO NGƯỜI PHỤ TRÁCH (RW)]:
% - Đoạn 1: Vấn đề thực tế lừa đảo viễn thông tại VN (thiệt hại hàng ngàn tỷ đồng, teencode lẩn tránh).
% - Đoạn 2: Các nghiên cứu trước đã làm gì (BERT, PhoBERT, LLMs).
% - Đoạn 3: GAP - 3 điểm nghẽn: loss đối xứng, alarm fatigue của LLM, độ trễ viễn thông.
% - Đoạn 4: 4 đóng góp cụ thể của bài báo (dạng itemize).
% - Đoạn 5: Bố cục các phần tiếp theo của bài báo.
%
% [TODO: Rà soát câu từ, bổ sung số liệu báo chí/chính phủ nếu muốn mở rộng Đoạn 1]

\section{Introduction}

% --- ĐOẠN 1: BỐI CẢNH THỰC TIỄN & NGUY CƠ VIỄN THÔNG TẠI VIỆT NAM ---
Digital telecommunications fraud, smishing (SMS phishing), and impersonation attacks have surged dramatically across Southeast Asia, inflicting estimated damages exceeding hundreds of millions of dollars annually in Vietnam alone~\cite{cam2026vnsed}. Attackers systematically exploit deceptive social engineering tactics, targeting vulnerable mobile users with fraudulent banking alerts, fake tax arrears notices, and fictitious job opportunities. In Vietnam, this operational landscape is characterized by extensive linguistic evasion tactics: scammers intentionally incorporate colloquial slang, compound teencode abbreviations, irregular spacing, and diacritic-stripped text to circumvent static keyword blacklists and carrier-level regex filters.

% --- ĐOẠN 2: CÁC TIẾP CẬN STATE-OF-THE-ART HIỆN TẠI ---
% [TODO: Có thể thêm 1-2 citation bài báo mới nếu muốn cập nhật xu hướng 2026]
Recent advances in natural language processing (NLP) have motivated researchers to deploy deep learning architectures and Large Language Models (LLMs) for message-based cyber defense~\cite{saias2025advances}. Standard transformer-based architectures, including PhoBERT~\cite{nguyen2020phobert} and multilingual variants~\cite{nguyenxuan2026mixed}, have demonstrated competitive overall classification accuracy on benchmark corpora. Simultaneously, prompt-based LLMs such as GPT-4 and Gemini have shown emergent reasoning capabilities in zero-shot and few-shot fraudulent intent classification~\cite{mahendru2024securenet}.

% --- ĐOẠN 3: KHOẢNG TRỐNG NGHIÊN CỨU (GAP STATEMENT) ---
% [LƯU Ý: Đoạn này phải khớp 100% với GAP trong §2 Related Work và claims.md]
However, in the empirical literature we reviewed, existing systems encounter a critical tripartite bottleneck. First, conventional neural classifiers are universally trained under symmetric loss functions (such as unweighted cross-entropy), treating false positives (misidentifying a legitimate message as scam) and catastrophic false negatives (failing to intercept a scam message that causes financial ruin) as mathematically identical. Second, state-of-the-art cloud LLMs suffer from severe ``alarm fatigue'': hyper-vigilant prompting causes precision to collapse as benign promotional and transactional notifications are mistakenly flagged~\cite{zhou2026fraudsmswalker}. Third, cloud API calls incur multi-second latencies ($3\text{--}6\text{ seconds/message}$) and recurring monetary costs, violating the strict sub-second latency and offline privacy constraints required for mobile edge deployment~\cite{seo2024ondevice, mrinal2026resilient}.

% --- ĐOẠN 4: 4 ĐÓNG GÓP CỐT LÕI CỦA BÀI BÁO (CONTRIBUTIONS) ---
To resolve these challenges, this paper presents an evidence-grounded empirical study and architectural solution. Our primary contributions are summarized as follows:
\begin{itemize}
    \item \textbf{Rigorous Multi-Seed Benchmark:} We establish a standardized benchmark on $2,665$ real-world Vietnamese SMS messages ($N=267$ isolated frozen test set) across 5 deterministic random seeds, comparing fine-tuned Vietnamese language models (PhoBERT-base, PhoBERT-large, FPT ViBERT) against few-shot cloud LLMs (Gemini 3.5 Flash).
    \item \textbf{Cost-Sensitive Social Media Modeling:} We fine-tune ViSoBERT~\cite{nguyen2023visobert} under a Cost-Sensitive Weighted Binary Cross-Entropy loss ($\mathcal{L}_{\text{WBCE}}$, $\alpha=5.0$), demonstrating a dominant Scam Recall of $95.62\% \pm 2.25\%$ and Macro-F1 of $92.51\% \pm 2.39\%$.
    \item \textbf{Statistical and Pareto Profiling:} Through McNemar's paired tests and deployment profiling, we establish that ViSoBERT attains Pareto-optimality over large-scale encoders ($3.8\times$ smaller than PhoBERT-large) while operating at $\approx 15\text{ ms/message}$ on consumer CPU hardware.
    \item \textbf{Two-Tier Architecture Grounding:} Based on the empirical breakdown of single-LLM precision collapse, we formulate an evidence-backed two-tier triage pipeline transitioning ambiguous edge samples to a specialized multi-agent forensic council.
\end{itemize}

% --- ĐOẠN 5: BỐ CỤC BÀI BÁO ---
The rest of this paper is organized as follows: Section~\ref{sec:related} reviews related work across three core research themes. Section~\ref{sec:method} details the dataset, loss formulation, and experimental pipeline. Section~\ref{sec:results} reports empirical results and statistical testing. Section~\ref{sec:discussion} discusses trade-offs and error analyses. Section~\ref{sec:threats} examines threats to validity, and Section~\ref{sec:conclusion} concludes the paper.
